% Lingua, codifica e impaginazione
\usepackage[italian]{babel}
\usepackage[T1]{fontenc}
\usepackage[utf8]{inputenc}
\usepackage{lmodern}
\usepackage[a4paper,margin=2.8cm]{geometry}
\usepackage{microtype}

% Matematica, codice e figure
\usepackage{amsmath,amssymb}
\usepackage{graphicx}
\usepackage{booktabs}
\usepackage{siunitx}
\usepackage{xcolor}
\usepackage{listings}
\usepackage{float}

% Collegamenti e indice
\usepackage[hidelinks]{hyperref}
\usepackage{bookmark}

% Ambienti per definizioni e osservazioni
\usepackage{amsthm}
\newtheorem{definizione}{Definizione}[chapter]
\newtheorem{teorema}{Teorema}[chapter]
\newtheorem{esempio}{Esempio}[chapter]
\theoremstyle{remark}
\newtheorem*{osservazione}{Osservazione}

% Stile del codice
\lstset{
	basicstyle=\ttfamily\small,
	frame=single,
	breaklines=true,
	columns=fullflexible,
	numbers=left,
	numberstyle=\tiny\color{gray},
	keywordstyle=\color{blue!70!black},
	commentstyle=\color{green!40!black},
	showstringspaces=false
}

\usepackage[T1]{fontenc}
\usepackage[utf8]{inputenc}
\usepackage[italian]{babel}

\theoremstyle{definition}
\newtheorem{definition}{Definizione}[chapter]
\newtheorem{example}{Esempio}[chapter]

\theoremstyle{remark}
\newtheorem{remark}{Osservazione}[chapter]
